\documentclass[conference,a4paper,10pt]{IEEEtran}
\IEEEoverridecommandlockouts

\usepackage[T1]{fontenc}
\usepackage[utf8]{inputenc}
\usepackage{lmodern}
\usepackage{microtype}
\usepackage{textcomp}
\usepackage{xcolor}
\usepackage{booktabs}
\usepackage{array}
\usepackage{tabularx}
\usepackage{listings}
\usepackage{url}
\usepackage[hyphens]{xurl}
\usepackage{hyperref}
\hypersetup{
  colorlinks=true,
  linkcolor=black,
  citecolor=black,
  urlcolor=black,
  pdftitle={Passive Security Analysis of Meshtastic Packet Captures},
  pdfauthor={Ali Murtaza Bhutto}
}

\lstset{
  basicstyle=\ttfamily\footnotesize,
  breaklines=true,
  columns=fullflexible,
  keepspaces=true,
  frame=single,
  framerule=0.2pt,
  framesep=2pt,
  xleftmargin=4pt,
}

\begin{document}

\title{Passive Security Analysis of Meshtastic Packet Captures}

\author{\IEEEauthorblockN{Ali Murtaza Bhutto}
\IEEEauthorblockA{ORCID: 0009-0007-2787-943X\\
\texttt{alibhutto101112@gmail.com}}}

\maketitle

\begin{abstract}
Meshtastic is a community-developed LoRa mesh networking stack that
ships in low-cost hardware and is widely deployed by hikers, ham
operators, and emergency-services volunteers. The protocol's defaults
trade strong security guarantees for ease of bootstrapping: every
device leaves the factory using one of seven publicly known modem
presets and a one-byte default pre-shared key index, with the source
and destination node identifiers transmitted in plaintext on every
frame. This paper describes \emph{meshtastic-security-audit}, an
open-source toolkit that consumes a PCAP capture of Meshtastic frames
and reports four classes of finding: use of the documented default
channel, node-identifier enumeration exposure, plaintext payloads,
and GPS coordinate leakage in unencrypted position frames. The
toolkit's parser, four detectors, reporter, and CLI run against a
synthetic Scapy-crafted corpus of five frames and recover 9 of 9
expected detections (precision and recall both 1.000). The artefact
is MIT-licensed and intended for authorised network audits.
\end{abstract}

\begin{IEEEkeywords}
Meshtastic, LoRa, LPWAN, mesh networking, security audit, packet
analysis, AES-CTR, GPS leakage, passive analysis
\end{IEEEkeywords}

\section{Introduction}

Low-power wide-area networks (LPWANs) have moved from industrial
telemetry into general-purpose human-to-human communication, driven
by hobbyist hardware (Heltec, T-Beam, RAK) and the open-source
Meshtastic firmware. A typical out-of-the-box deployment routes
text, position, and telemetry messages across kilometres of terrain
at sub-watt transmit power, with no infrastructure beyond the nodes
themselves. The convenience of that model rests on a small set of
defaults: a fixed list of seven modem presets, a documented one-byte
PSK index for the primary channel, and an unencrypted header that
identifies the source and destination of every frame.

Each of those defaults is well known and intentional---the
Meshtastic project documents them and they are necessary for
interoperability between strangers in an ad-hoc mesh. They are also
not what an operator running a sensitive deployment would choose if
they understood the consequences. Two patterns recur in the
practitioner discussions: an operator forgets to rotate the PSK on
deployment, leaving traffic decodable by anyone with a Meshtastic
node; and an operator enables position broadcasts on an unencrypted
channel, broadcasting their location to every receiver in range.

This paper introduces a passive analysis toolkit that examines a
PCAP capture of Meshtastic frames and reports the four most common
classes of exposure. The toolkit is not a decryption attack on
correctly configured devices and does not transmit, inject, or
otherwise touch a live mesh.

\paragraph{Contribution.}
(1) A Scapy-based parser that recovers the Meshtastic frame structure
from a PCAP irrespective of the link-layer encapsulation used at
capture time. (2) Four independent detectors, each emitting findings
in a severity band that is consistent across the toolkit. (3) An HTML
and JSON reporter that renders the findings for both humans and
downstream tooling. (4) A synthetic, byte-level test corpus and an
eval that measures the detectors against ground truth without
requiring RF hardware.

\section{Background}

\paragraph{Meshtastic frame structure.}
A Meshtastic frame as it travels over the LoRa physical layer
\cite{lora_phy} consists of a 14- or 16-byte header (depending on
firmware revision) followed by a length-variable payload. The header
fields the toolkit consumes are: a 32-bit destination node
identifier, a 32-bit source node identifier, a 32-bit packet
identifier, a flags byte (carrying a 3-bit hop limit, a want-ack
bit, and an encrypted-payload bit), and an 8-bit channel hash. The
channel hash is the XOR of every byte of the channel name with every
byte of the PSK \cite{meshtastic_docs}.

\paragraph{Encryption.}
When the encrypted-payload bit is set, the payload is encrypted with
AES-256 in counter mode (CTR), using the packet identifier and source
node id as the nonce material. CTR \cite{nist80038a} is appropriate
for short messages but provides no authentication; a separate
authenticated mode would be preferable for a security-sensitive
deployment, though that would break compatibility with the existing
fleet.

\paragraph{Default channel.}
For the primary channel, the firmware uses a one-byte PSK index of
\texttt{0x01}, which is expanded by the firmware into a 256-bit AES
key. Any device using the same preset name and the same default
index can decrypt the traffic. The toolkit detects this by computing
the channel hash that would result for each documented preset name
under the default index and checking whether the captured frames
carry that hash.

\paragraph{Threats to confidentiality and location privacy.}
The combination of plaintext source identifiers and unencrypted
position payloads is sufficient, on its own, to track a node across
days of capture. The LPWAN security literature
\cite{adelantado2017,enisa2024} has flagged these classes of
exposure as recurring failure modes across LoRaWAN, Sigfox, and
Meshtastic deployments.

\section{Toolkit architecture}

\subsection{Parser}

\texttt{audit/pcap\_parser.py} reads a PCAP via Scapy and, for each
record, attempts to interpret the raw bytes as either the current
16-byte Meshtastic header or the legacy 14-byte variant. A small
plausibility filter (non-zero source id, non-zero packet id, sane
broadcast/unicast destination) screens out non-Meshtastic frames
without raising. The result is a list of \texttt{MeshFrame}
dataclasses with the header fields broken out and the encrypted
payload preserved verbatim.

\subsection{Detectors}

Four detectors run independently against the parsed frames; each
emits one or more \texttt{Finding} objects.

\begin{itemize}
  \item \textbf{\texttt{default\_key}} computes the channel hash for
    every documented preset name under the default PSK index and
    flags any frame whose channel byte matches. Severity:
    \texttt{high}.
  \item \textbf{\texttt{node\_enum}} counts distinct source and
    destination node ids across the capture, emitting an \texttt{info}
    finding with the source inventory and a \texttt{medium} finding
    when more than one unicast destination is addressed.
  \item \textbf{\texttt{unencrypted}} reports two complementary
    signals: the encrypted bit being clear in the flags byte (high),
    and a payload whose first bytes parse as a plaintext Meshtastic
    \texttt{Data} protobuf (\texttt{critical}). The protobuf check is
    a stronger signal than the flag bit because it is content-based.
  \item \textbf{\texttt{gps\_leak}} decodes any \texttt{Data}
    envelope on port 3 (\texttt{POSITION\_APP}) and extracts the
    \texttt{latitude\_i} and \texttt{longitude\_i} fields. When a
    valid lat/lon pair is recovered, it emits a \texttt{critical}
    finding. The detector never attempts to decrypt encrypted
    payloads; an encrypted position frame produces no finding.
\end{itemize}

\subsection{Reporter and CLI}

\texttt{audit/reporter.py} renders the finding set as either an
inline-styled HTML page (no external assets) or a JSON document with
a stable summary block. The CLI (\texttt{cli.py}, Click-based)
exposes \texttt{--pcap}, \texttt{--output}, \texttt{--format},
\texttt{--detector} (repeatable), and \texttt{--severity-threshold}.

\section{Evaluation}

A live regression run on the synthetic corpus in
\texttt{tests/fixtures/test\_packets.py} produces the results in
Table~\ref{tab:eval}. The corpus contains five frames covering all
four detector signals; ground-truth labels live in
\texttt{eval/labelled\_corpus.json} and the harness lives in
\texttt{eval/run\_eval.py}.

\begin{table}[!t]
\caption{Per-detector confusion-matrix metrics on the synthetic
labelled corpus (5 frames). The corpus is intentionally small and
controlled; per-detector ground truth is the union of detector
signals each frame is constructed to exhibit. Numbers reproduce by
running \texttt{python eval/run\_eval.py}.}
\label{tab:eval}
\centering
\footnotesize
\begin{tabularx}{\columnwidth}{Xrrrrrr}
\toprule
\textbf{Detector} & \textbf{TP} & \textbf{FP} & \textbf{FN} &
\textbf{P} & \textbf{R} & \textbf{F$_1$} \\
\midrule
\texttt{default\_key}  & 1 & 0 & 0 & 1.000 & 1.000 & 1.000 \\
\texttt{node\_enum}    & 5 & 0 & 0 & 1.000 & 1.000 & 1.000 \\
\texttt{unencrypted}   & 2 & 0 & 0 & 1.000 & 1.000 & 1.000 \\
\texttt{gps\_leak}     & 1 & 0 & 0 & 1.000 & 1.000 & 1.000 \\
\midrule
\textbf{Aggregate}     & 9 & 0 & 0 & 1.000 & 1.000 & 1.000 \\
\bottomrule
\end{tabularx}
\end{table}

The pytest suite (\texttt{tests/}) carries 23 tests: 7 for the
parser, 8 for the detectors, 4 for the reporter (HTML, JSON, severity
ordering, empty-input handling), and 4 end-to-end CLI tests via the
Click \texttt{CliRunner}. The full suite executes in approximately
1.05 seconds on the reference host.

\subsection{DevSecOps gates}

The release passes three gates:

\begin{itemize}
  \item \textbf{Bandit} (medium severity and above) reports zero
    findings across 840 lines of Python.
  \item \textbf{pip-audit} reports no known vulnerabilities in the
    installed dependency closure after pinning \texttt{pytest$\geq$9.0.3}
    and upgrading \texttt{pip}.
  \item \textbf{Semgrep} with the \texttt{p/python} and
    \texttt{p/security-audit} rule packs exits 0.
\end{itemize}

Raw scan outputs are committed under \texttt{results/security\_scan.md}.

\section{Threat model}

The toolkit is read-only. It produces no RF traffic, opens no network
sockets, and does not attempt to decrypt encrypted payloads. The
attack surface is therefore the PCAP parser itself: a malformed or
adversarial PCAP could crash the parser, but the plausibility filter
and Scapy's own input handling bound the impact to a Python
exception. The CLI accepts only local file paths.

The toolkit is intentionally narrow about what it does \emph{not}
attempt:

\begin{itemize}
  \item No active RF capture. The operator must produce the PCAP via
    an SDR pipeline or the firmware's debug serial sniffer.
  \item No decryption attack on correctly configured channels. A
    user who has rotated the PSK is outside the toolkit's reach;
    that is the intended outcome.
  \item No node identity resolution. The output stays at the
    \texttt{!xxxxxxxx} layer; mapping that identifier to a real
    human is an operator decision and is governed by the
    authorisation chain.
\end{itemize}

\section{Conclusion}

\emph{meshtastic-security-audit} is a small, focused toolkit that
operationalises four well-understood Meshtastic-specific exposures
into a single command-line audit. The contribution is not a new
attack; the defaults this toolkit flags have been documented by the
Meshtastic project for years \cite{meshtastic_docs} and are
unavoidable consequences of the protocol's bootstrapping model. The
contribution is reproducible, measured, low-effort auditing of the
protocol surface, so an operator deploying Meshtastic in a
sensitive environment can answer "what would an observer with a
receiver in range see?" before deploying rather than after.

\begin{thebibliography}{9}

\bibitem{meshtastic_docs}
Meshtastic Project, ``Meshtastic protocol documentation,'' 2024.
\url{https://meshtastic.org/docs/overview/} and the protobuf
definitions at \url{https://github.com/meshtastic/protobufs}.

\bibitem{lora_phy}
Semtech Corporation, ``LoRa modulation basics,'' Application Note
AN1200.22, 2015.
\url{https://www.semtech.com/products/wireless-rf/lora-connect}

\bibitem{nist80038a}
M.~Dworkin, ``Recommendation for block cipher modes of operation:
Methods and techniques,'' National Institute of Standards and
Technology, NIST Special Publication 800-38A, Dec.\ 2001.

\bibitem{adelantado2017}
F.~Adelantado, X.~Vilajosana, P.~Tuset-Peiro, B.~Martinez,
J.~Melia-Segui, and T.~Watteyne, ``Understanding the limits of
LoRaWAN,'' \emph{IEEE Communications Magazine}, vol.~55, no.~9,
pp.~34--40, Sep.\ 2017.

\bibitem{enisa2024}
European Union Agency for Cybersecurity (ENISA), ``ENISA threat
landscape 2024,'' Sep.\ 2024.
\url{https://www.enisa.europa.eu/publications/enisa-threat-landscape-2024}

\end{thebibliography}

\end{document}
